\chapter{More on Functions in Python: Parameters, Return Values, and Scope}

\section{Introduction: Real-Life Hook}

Think about daily routines:

\begin{itemize}
    \item Making tea: you can vary sugar or milk each time  
    \item Calculating exam scores for different subjects  
    \item Sending personalized greetings to friends  
\end{itemize}

In Python, \textbf{functions can take inputs (parameters)}, \textbf{return outputs}, and have \textbf{local vs global variables} to make programs flexible and reusable.

\section{Concept: Function Parameters and Arguments}

\begin{definition}{Function Parameter}
A \textbf{parameter} is a variable listed in a function definition to receive input when the function is called.
\end{definition}

\subsection{Example: Personalized Greeting}

\begin{verbatim}
def greet(name):
    print("Hello,", name)

greet("Amina")
greet("Rafi")
\end{verbatim}

- \texttt{name} is the \textbf{parameter}  
- \texttt{"Amina"} and \texttt{"Rafi"} are \textbf{arguments}

\section{Default Parameters}

\begin{verbatim}
def greet(name="Friend"):
    print("Hello,", name)

greet()         # Hello, Friend
greet("Amina")  # Hello, Amina
\end{verbatim}

- Functions can have default values if arguments are not provided

\section{Return Values}

\begin{verbatim}
def add(a, b):
    return a + b

result = add(5, 3)
print("Sum is:", result)
\end{verbatim}

- \texttt{return} sends a value back to the caller  
- Returned values can be stored, printed, or used in calculations

\section{Variable Scope: Local vs Global}

\begin{verbatim}
x = 10  # global variable

def change_x():
    x = 5  # local variable
    print("Inside function:", x)

change_x()        # Inside function: 5
print("Outside function:", x)  # 10
\end{verbatim}

- Variables inside a function are \textbf{local}  
- Variables outside are \textbf{global}  
- Understanding scope prevents accidental errors

\section{Docstrings: Documenting Functions}

\begin{verbatim}
def add(a, b):
    """Returns the sum of a and b"""
    return a + b

print(add.__doc__)
\end{verbatim}

- Docstrings describe what a function does  
- Helpful for readability and teaching others

\section{Mini-Project: Temperature Converter}

Create a program with a function to convert Celsius to Fahrenheit:

\begin{verbatim}
def celsius_to_fahrenheit(c):
    """Converts Celsius to Fahrenheit"""
    return (c * 9/5) + 32

temp_c = float(input("Enter temperature in Celsius: "))
print("Temperature in Fahrenheit:", celsius_to_fahrenheit(temp_c))
\end{verbatim}

\section{Step-by-Step Diagram: Function with Parameters and Return Values}

\begin{center}
\begin{tikzpicture}[node distance=1.5cm]
\node (start) [draw, ellipse, fill=green!20] {Start Program};
\node (define) [draw, rectangle, fill=yellow!20, below of=start] {Define Function with Parameters};
\node (call) [draw, rectangle, fill=orange!20, below of=define] {Call Function with Arguments};
\node (execute) [draw, rectangle, fill=orange!20, below of=call] {Execute Function Body};
\node (return) [draw, rectangle, fill=yellow!20, below of=execute] {Return Value};
\node (use) [draw, rectangle, fill=orange!20, below of=return] {Use Value in Program};
\node (end) [draw, ellipse, fill=green!20, below of=use] {End Program};
\draw[->] (start) -- (define);
\draw[->] (define) -- (call);
\draw[->] (call) -- (execute);
\draw[->] (execute) -- (return);
\draw[->] (return) -- (use);
\draw[->] (use) -- (end);
\end{tikzpicture}
\end{center}

\section{Key Concepts}

\begin{itemize}
    \item Functions can take \textbf{parameters} and receive \textbf{arguments}  
    \item Functions can \textbf{return values} to be used elsewhere  
    \item \textbf{Local vs global variables}: scope matters  
    \item \textbf{Docstrings} document what a function does for readability  
\end{itemize}

\section{End-of-Chapter Exercises}

\begin{exercise}{Very Short Questions}
\begin{enumerate}
    \item What is a function parameter?  
    \item How do you return a value from a function?  
    \item What is the difference between local and global variables?  
\end{enumerate}
\end{exercise}

\begin{exercise}{Short Answer Questions}
\begin{enumerate}
    \item Write a function that takes two numbers and returns their difference.  
    \item Explain why function docstrings are useful.  
    \item What happens if you change a global variable inside a function without the global keyword?  
\end{enumerate}
\end{exercise}

\begin{exercise}{Activity-Based Questions}
\begin{enumerate}
    \item Create a program with a function to convert kilometers to miles and test it with user input.  
    \item Write a function that takes a list of numbers and returns the maximum value.  
    \item Make a program that uses a function to greet a user with their name and age. Include a docstring.  
\end{enumerate}
\end{exercise}

\section{Teacher Notes}

\marginnote{
\textbf{Learning Objectives:}  
\begin{itemize}
    \item Students understand parameters, arguments, return values, and variable scope  
    \item Students can create reusable and documented functions  
    \item Students can organize programs efficiently  
\end{itemize}

\textbf{Teaching Tips:}  
\begin{itemize}
    \item Start with simple functions and gradually introduce parameters and return values  
    \item Demonstrate local vs global variables using clear examples  
    \item Encourage writing docstrings for all functions  
\end{itemize}

\textbf{Assessment:}  
\begin{itemize}
    \item Check correctness of functions with parameters and return values  
    \item Observe understanding of scope and docstrings  
\end{itemize}
}